\emph{Setting.} A tower representative is a signature $(p,q)$; its residue is $p - q \bmod 8$, which the $(1,1)$ ascent preserves (Fact \claimref{dark-matter.imported-eight-towers-mod-8-clock}). The check builds each algebra with \texttt{GA.py}, the representative's own generators first and the eight added ones last.

(i) Let $s = (p,q)$ be causal of dimension $n$, with majority sign $\sigma$ ($\sigma = +1$ if $p > q$), and let $g_{1}, \dots, g_{8}$ be the added vectors of $\beta s$, all of square $\sigma$. Their product $\omega_{8}$ has $\omega_{8}^{2} = (-1)^{28}\sigma^{8} = 1$. Each old generator $e_{i}$ anticommutes with every $g_{j}$, hence commutes with $\omega_{8}$ (eight sign changes), while each $g_{j}$ anticommutes with $\omega_{8}$ (it meets seven others). So $a_{i} = e_{i}\omega_{8}$ has $a_{i}^{2} = e_{i}^{2}\omega_{8}^{2} = e_{i}^{2}$, $a_{i}a_{k} = e_{i}e_{k}$ for $i \ne k$ (so the $a_{i}$ anticommute), and $a_{i}g_{j} = e_{i}\omega_{8}g_{j} = -e_{i}g_{j}\omega_{8} = g_{j}e_{i}\omega_{8} = g_{j}a_{i}$. Reversion fixes $a_{i}$, a blade of grade $9$ ($(-1)^{36} = 1$), and each $g_{j}$. By the universal property the relations give an algebra map $\Cl(p,q) \to A = \langle a_{i}\rangle$; the products $a_{S}$ over subsets $S$ are $\pm e_{S}$ or $\pm e_{S}\omega_{8}$, distinct blades, so it is injective. The products $a_{S}g_{T}$ over all subsets are $2^{n+8}$ distinct blades up to sign (computed for $(2,1)$ and $(4,1)$; the mirror cases differ only in the signs of the squares), so $A \otimes B \to \Cl(\beta s)$ is onto, and by dimension an isomorphism; reversion of $\Cl(\beta s)$ restricts to the reversion of $A$ (it fixes the generators $a_{i}$ and reverses products) and of $B$. $B$ is $\Cl(8,0)$ or $\Cl(0,8)$, of residue $0$, a matrix algebra over $\R$ by the imported classification (Fact \claimref{dark-tower-matter.imported-clifford-structure}), of dimension $256$, so $B \cong \Mat(16,\R)$. \emph{Computed:} with $h = g_{1234}, g_{5678}, g_{1256}, g_{1357}$ (pairwise overlaps even, so they commute; each squares to $1$ for either sign), $f = \prod\tfrac12(1 + h)$ has $f^{2} = f$, $\widetilde{f} = f$, scalar part $\tfrac1{16}$, and $fxf \in \R f$ for all $256$ blades $x$, nonzero for exactly sixteen; so $fBf = \R f$ and $f$ is primitive. For $\psi, \varphi \in Bf$, $\widetilde{\psi}\varphi = c(\psi,\varphi)f$, and reversing gives $c(\varphi,\psi)f = c(\psi,\varphi)\widetilde{f} = c(\psi,\varphi)f$: the reversion pairing on $Bf \cong \R^{16}$ is symmetric. Therefore $\Cl(\beta s) \cong A \otimes \Mat(16,\R) \cong \Mat(16, \Cl(s))$, Morita equivalent to $\Cl(s)$ (Fact \claimref{tower.morita-equivalence-algebra-amplification}); its minimal left ideals are $S \otimes Bf$ with $S$ one of $A$, and its reversion pairing is that of $S$ tensored with a symmetric form, of the same symmetry type. The type data: $\beta$ adds $8$ to $p - q$ or subtracts it, so the residue is unchanged; $\min$ stays $1$; $n$ changes by $8$; the pseudoscalar square $(-1)^{n(n-1)/2}(-1)^{q}$ is multiplied by $(-1)^{8n + 28}$, and by $(-1)^{8}$ when $q$ grows, both equal to $1$ (computed for the four odd representatives); for odd $n$ the pseudoscalar is central in both (computed). Reversal $(p,q) \mapsto (q,p)$ commutes with $\beta$, since it exchanges the majority sign.

(ii) Premise by premise, at the pinned versions. \emph{F1 to F4.} The foundation's directions with sign bits generate $\Cl(p,q)$ for every signature; the check confirms the Clifford relations of every representative of both models. \emph{The clock.} Every representative of both models has the residue it represents (computed). \emph{The single-clock filter.} All populated representatives are causal in both (computed). \emph{P1} (\claimref{dark-matter.p1-single-frozen-point}). It freezes one $I$ and couples every tower to one metric; it places no condition on a tower's signature (as \claimref{synthesis.frozen-point-complement-window-lowest-causal} records, nothing in the series says the towers' spacetimes lie in the frozen point's complement). Its orientation clause acts on the reversal orbit $\{2, 6\}$, whose representatives $(3,1)$, $(1,3)$ are the same in both models and have dimension $0 \bmod 4$, so both carry the flag (Theorem \claimref{synthesis.oriented-causal-towers-mirror-pair-every-window}: oriented exactly when $n \equiv 0 \pmod 4$); the odd representatives have dimensions $3, 5$ and $11, 13$, odd in both, so no odd tower is oriented in either. \emph{P2} (\claimref{dark-matter.p2-indifference}). It equates towers only along $r \mapsto -r$ and reaches labels; each model assigns to $-r$ the reversal of its signature for $r$ (computed), and the labels (residues and tower charges) are the same in both. \emph{The costed clause.} Residues $0$ and $4$ are treated alike in both models. \emph{H1 and the constituent definitions.} They are stated for a spinor ideal of any signature; by (i) the ideal of $\beta s$ is that of $s$ tensored with $\R^{16}$ with a symmetric reversion form, so every pairing statement reads the same, and the momentum map is defined in both. \emph{Anomaly import.} Local gauge and gravitational anomaly forms have degree $n + 2$ and exist only in even dimension (and the gravitational one only for $n \equiv 2 \pmod 4$, Theorem \claimref{synthesis.residues-zero-four-real-chirality-anomaly}(i), (ii)); the odd representatives have odd $n$ in both models. The anomaly obstruction of that theorem's clause (iii) concerns residue $0$ at $(1,1)$, which the costed clause treats alike in both. \emph{Inverse-square law.} It bounds an extra direction carrying the one metric; in both models the towers' unshared directions are their own (Proposition \claimref{dark-matter.thickness-not-gravitational-quantity}), so it is met alike. \emph{R1 to R3.} Their dimension-generic forms hold in both; R2's parenthetical list of dimensions $2$, $3$, $4$ is the window P3 supplies and is read as such. \emph{Remaining premises and imports} (the quasi-free premise (Q), tameness, the indifference premise, the frozen-point, tractor and field-theory imports) mention no tower's signature and are met or not alike. \emph{Selectors.} The flag bit reads $n \bmod 4$, real chirality reads $n \bmod 4$ and the pseudoscalar, the anomaly degree reads $n$'s parity and $n \bmod 4$, the pairing type reads the reversion form: all unchanged by (i). So no premise other than P3 separates the models; iterating $\beta$ gives a model for every $k \ge 0$ with the odd towers at dimensions $3 + 8k$ and $5 + 8k$.

(iii) Under (W) with $\Lambda > 0$ a causal signature is admitted exactly when $\max(p,q) \le 4$ (\claimref{synthesis.frozen-point-complement-window-lowest-causal}(ii), (iii)). Every representative of $M_{L}$ passes; $(10,1)$, $(1,10)$, $(12,1)$, $(1,12)$ fail. An enumeration of all causal signatures of odd residue up to dimension sixteen finds exactly one admitted member per residue, the lowest. (W) is a selector that is not periodic along a residue, which (ii) shows the descent lacks.

(iv) With $V_{2k} = \pi^{k}/k!$ and $V_{2k+1} = 2^{k+1}\pi^{k}/(2k+1)!!$, the shares $V_{3}/V_{D}$ are $4/3$ ($D = 2$), $8/3\pi$ ($D = 4$), $160/\pi^{4}$ ($D = 10$) and $960/\pi^{5}$ ($D = 12$), with the mirror's share $1$ in both models. So the ratio is $1 + 8/3 + 16/3\pi = (11\pi + 16)/3\pi$ in $M_{L}$ (Proposition \claimref{dark-matter.unit-ball-rule}) and $1 + 320/\pi^{4} + 1920/\pi^{5}$ in $M_{8}$. The check brackets both with $\tfrac{333}{106} < \pi < \tfrac{355}{113}$: $(5.36, 5.37)$ and $(10.5, 10.6)$, disjoint. Since the observed ratio is the only datum separating the models, the window is an input that the fit selects, which is the status Observation \claimref{dark-matter.status-rule} records for the relation; P3 remains a premise, now with its lowest-causal clause shown independent of the descent.
